% !TeX program = xelatex
% 临时交互样例；从项目 Book1 目录编译，不改教材或共享样式。
\documentclass[UTF8,zihao=-4,openany,fontset=fandol]{ctexbook}
\newcommand{\BookTitle}{解答隐藏交互样例}
\newcommand{\BookVolumeName}{临时试样}
\newcommand{\BookNumber}{1}
\newcommand{\BookDraftDate}{2026年9月}
\input{../Shared/Preamble}
\usepackage{ocgx2}
\hypersetup{pdftitle={解答隐藏交互样例},pdfsubject={悬浮提示与点击展开},pdfhighlight=/N}
\newsavebox{\DemoAnswerBox}
\newlength{\DemoAnswerHeight}
\newlength{\DemoAnswerWidth}

% A: 答案，初始隐藏；B: 尚未展开，初始可见；H: 鼠标悬浮状态。
% 鼠标移入 H=on，移出 H=off；点击 A=on，B=H=off。
% 提示只在 B 与 H 都开启时出现，按钮附着在 B 上，展开后也会消失。
\NewDocumentCommand{\DemoHiddenAnswer}{m +m}{%
  \par\smallskip
  \setlength{\DemoAnswerWidth}{\linewidth}%
  \sbox{\DemoAnswerBox}{%
    \begin{minipage}{\DemoAnswerWidth}%
      \begin{solution}#2\end{solution}%
    \end{minipage}%
  }%
  \setlength{\DemoAnswerHeight}{\dimexpr\ht\DemoAnswerBox+\dp\DemoAnswerBox\relax}%
  \begin{ocg}[showingui=false,printocg=never]{#1 hidden}{#1-B}{1}\end{ocg}%
  \begin{ocg}[showingui=false,printocg=never]{#1 hover}{#1-H}{0}\end{ocg}%
  \noindent
  \begin{tikzpicture}[baseline=(answer.south)]
    \node[anchor=south west,inner sep=0pt,outer sep=0pt] (answer) at (0,0)
      {\begin{ocg}[printocg=ifvisible]{#1 solution}{#1-A}{0}\usebox{\DemoAnswerBox}\end{ocg}};
    \node[anchor=center,inner sep=0pt,outer sep=0pt] at (answer.center)
      {\begin{ocmd}{\AllOn{#1-B,#1-H}}%
        {\kaishu\fontsize{10}{12}\selectfont\color{AlgebraGray}点击显示解答}%
       \end{ocmd}};
    \node[anchor=south west,inner sep=0pt,outer sep=0pt] at (answer.south west)
      {\begin{ocg}[showingui=false,printocg=never]{#1 hidden}{#1-B}{1}%
        \actionsocg[onmouseall]{}
          {{#1-H},{},{},{#1-A}}
          {{},{#1-H},{},{#1-B,#1-H}}
          {\rule{0pt}{\DemoAnswerHeight}\rule{\DemoAnswerWidth}{0pt}}%
       \end{ocg}};
  \end{tikzpicture}
  \par\smallskip
}

\begin{document}
\pagestyle{plain}
\setcounter{chapter}{10}
\section*{解答隐藏交互样例}

请用桌面版 Adobe Acrobat Reader 打开。把鼠标移入题目下方留白区，中央会出现提示；点击该区域后，解答保持显示。页末可重新隐藏，便于反复试用。

\ProseParagraphBreak
两道题的解答独立控制，隐藏时保留原来的排版高度。

\begin{exercise}
求全部满足 \(x\equiv3\pmod8\)、\(x\equiv7\pmod{12}\) 的整数，并判断把第二个余数改为 \(5\) 后是否有解。
\end{exercise}
\DemoHiddenAnswer{one}{
两模数的最大公因子为 \(4\)，且 \(3\equiv7\pmod4\)，所以同余方程有解。写 \(x=3+8t\)，代入第二式得
\[
8t\equiv4\pmod{12},\qquad 2t\equiv1\pmod3,
\]
故 \(t\equiv2\pmod3\)。全部解为 \(x=19+24u\)，其中 \(u\in\mathbb Z\)。

若把第二个余数改为 \(5\)，则两个同余式会要求 \(x\) 模四同时余三与余一，矛盾，因此无解。
}

\begin{exercise}
设 \(R\) 为交换环，理想 \(I\) 中每个元素都是幂零元。证明 \(a\in R\) 可逆，当且仅当 \(a+I\) 在 \(R/I\) 中可逆。
\end{exercise}
\DemoHiddenAnswer{two}{
单位在环同态下仍为单位，因此必要性成立。

反过来，设 \(a+I\) 可逆，选逆元的代表 \(b\in R\)，则 \(ab=1+n\)，其中 \(n\in I\)。取正整数 \(m\) 使 \(n^m=0\)。有限几何级数给出
\[
(1+n)\Bigl(1-n+n^2-\cdots+(-n)^{m-1}\Bigr)=1.
\]
因此 \(1+n\) 是单位，\(a\) 的逆元为
\[
b\Bigl(1-n+n^2-\cdots+(-n)^{m-1}\Bigr).
\]
这里环交换，所以该元素同时是左逆与右逆，结论成立。
}

\vspace{1em}
\noindent{\small\color{AlgebraBlue}
\actionsocg{}{one-B,two-B}{one-A,two-A,one-H,two-H}{重新隐藏全部解答}
\qquad
\actionsocg{}{one-A,two-A}{one-B,two-B,one-H,two-H}{显示全部解答}}

\vspace{.5em}
\noindent{\footnotesize\color{AlgebraGray}临时效果样例。教材正文与正式PDF未作改动。}
\end{document}
